% RASCUNHO EXPERIMENTAL DE ARTIGO - ABNT. Requer validação da instituição.
% Base: classe abntex2, não representa certificação das normas mais recentes.
\documentclass[article,12pt,oneside,a4paper,brazil]{abntex2}
\usepackage[T1]{fontenc}
\usepackage[utf8]{inputenc}
\usepackage[brazil]{babel}
\usepackage{lmodern}
\usepackage{graphicx,amsmath,amssymb,booktabs}
\usepackage{hyperref,url}
\usepackage[alf]{abntex2cite}
% Compatibilidade com comandos de citação do motor PaperOrchestra.
\providecommand{\citep}[1]{\cite{#1}}
\providecommand{\citet}[1]{\citeonline{#1}}
\providecommand{\citealp}[1]{\cite{#1}}
\providecommand{\citealt}[1]{\citeonline{#1}}
\providecommand{\citeauthor}[1]{\citeonline{#1}}
\providecommand{\citeyear}[1]{\cite{#1}}
\titulo{Título provisório do artigo científico}
\autor{Autor(a) — preencher antes da publicação}
\local{Brasil}
\data{\the\year}
\begin{document}
\selectlanguage{brazil}
\maketitle
\begin{resumo}
%%%%%%%%% RESUMO %%%%%%%%%
Esta é uma estrutura preliminar. O resumo deve descrever o objetivo,
o método e as conclusões sustentadas pelos dados. Não inventar resultados.
%%%%%%%%% RESUMO %%%%%%%%%
\end{resumo}
\textual
\section{Introdução}
\label{sec:intro}
%%%%%%%%% INTRODUÇÃO %%%%%%%%%
Apresente o problema e a justificativa.
%%%%%%%%% INTRODUÇÃO %%%%%%%%%
\section{Referencial teórico}
\label{sec:related_work}
%%%%%%%%% REVISÃO %%%%%%%%%
Discuta a literatura com referências verificadas.
%%%%%%%%% REVISÃO %%%%%%%%%
\section{Metodologia}
\label{sec:method}
%%%%%%%%% MÉTODO %%%%%%%%%
Descreva o método realmente utilizado.
%%%%%%%%% MÉTODO %%%%%%%%%
\section{Resultados e discussão}
\label{sec:experiments}
%%%%%%%%% RESULTADOS %%%%%%%%%
Se não há dados, declare explicitamente que não houve experimentos.
%%%%%%%%% RESULTADOS %%%%%%%%%
\section{Considerações finais}
\label{sec:conclusion}
%%%%%%%%% CONCLUSÃO %%%%%%%%%
Indique conclusões sustentadas pelos materiais.
%%%%%%%%% CONCLUSÃO %%%%%%%%%
\postextual
\bibliography{references}
\end{document}
